
\bibitem{bernstein1912}
S.~N.~Bernstein.
\emph{D\'emonstration du th\'eor\`eme de Weierstrass fond\'ee sur le
calcul des probabilit\'es}.
Communications de la Soci\'et\'e Math\'ematique de Kharkov (2) 13
(1912), 1--2.

\bibitem{cargoshisha1966}
G.~T.~Cargo and O.~Shisha.
\emph{The Bernstein form of a polynomial}.
Journal of Research of the National Bureau of Standards 70B (1966),
79--81. (Bernstein coefficients enclose the range on the simplex ---
the sound direction this calculus's kernel uses.)

\bibitem{farouki2012}
R.~T.~Farouki.
\emph{The Bernstein polynomial basis: a centennial retrospective}.
Computer Aided Geometric Design 29 (2012), no.~6, 379--419.

\bibitem{decasteljau1959}
P.~de~Casteljau.
\emph{Outillages m\'ethodes calcul}.
Technical report, A.~Citro\"en, Paris, 1959. (The founding
subdivision/evaluation algorithm for polynomials in Bernstein form.)

\bibitem{laneriesenfeld1980}
J.~M.~Lane and R.~F.~Riesenfeld.
\emph{A theoretical development for the computer generation and display
of piecewise polynomial surfaces}.
IEEE Transactions on Pattern Analysis and Machine Intelligence 2
(1980), no.~1, 35--46.

\bibitem{garloff1986}
J.~Garloff.
\emph{Convergent bounds for the range of multivariate polynomials}.
In: Interval Mathematics 1985 (K.~Nickel, ed.), Lecture Notes in
Computer Science 212, Springer, 1986, 37--56. (Subdivision sharpens
Bernstein enclosures --- the completeness partner of the enclosure
soundness used here.)
